% Part: turing-machines
% Chapter: undecidability
% Section: enumerating-tms

\documentclass[../../../include/open-logic-section]{subfiles}

\begin{document}

\olfileid{tur}{und}{enu}
\olsection{ट्यूरिंग यंत्रांचे प्रगणन}

\begin{explain}
सर्व ट्यूरिंग यंत्रांचा संच !!{enumerable} म्हणजे गणनीय
आहे, असे आपण दाखवू शकतो. प्रत्येक ट्यूरिंग यंत्राचे
मर्यादित वर्णन करता येते, ही त्यामागची बाब आहे. अवस्थांचा
संच आणि फितीवरील चिन्हसंच हे दोन्ही सांत संच आहेत.
संक्रमण फलन हे $Q \times \Sigma$ पासून
$Q \times \Sigma \times \{\TMleft, \TMright, \TMstay\}$
मधील आंशिक फलन आहे. म्हणून ते परिभाषित असलेल्या
मर्यादित संख्येच्या आदान-जोड्यांवरील त्याची मूल्ये
यादीत दिल्यास त्याचेही वर्णन करता येते.

ही मांडणी एका मर्यादेपर्यंत खरी असली, तरी त्यात
एक सूक्ष्म फरक आहे. ट्यूरिंग यंत्राच्या व्याख्येत
अवस्थांच्या संचात आणि फितीच्या वर्णमालेत कोणते
!!{element}s असावेत यावर बंधन घातलेले नाही.
म्हणून उदाहरणार्थ, प्रत्येक वास्तव संख्येला अवस्था
म्हणून वापरणारे ट्यूरिंग यंत्र तांत्रिकदृष्ट्या
असू शकते. मात्र ट्यूरिंग यंत्राचे \emph{वर्तन}
कोणत्या वस्तू अवस्था किंवा फितीवरील चिन्हे म्हणून
वापरल्या जातात यावर अवलंबून नसते.
\olref{fig:variants} मधील दोन ट्यूरिंग यंत्रे पाहा.
\begin{figure}
\begin{center}
\begin{tikzpicture}[->,>=stealth',shorten >=1pt,auto,node distance=2.8cm,
                    semithick]
  \tikzstyle{every state}=[fill=none,draw=black,text=black]

  \node[initial,state]         (A)              {$q_0$};
  \node[state]         (B) [right of=A] {$q_1$};

  \path (A) edge [bend left] node {\TMtrans{\TMstroke}{\TMstroke}{\TMright}} (B)
        (B) edge [loop above] node {\TMtrans{\TMblank}{\TMblank}{\TMright}} (B)
            edge [bend left] node {\TMtrans{\TMstroke}{\TMstroke}{\TMright}} (A);
\end{tikzpicture}\\
\begin{tikzpicture}[->,>=stealth',shorten >=1pt,auto,node distance=2.8cm,
  semithick]
\tikzstyle{every state}=[fill=none,draw=black,text=black]

\node[initial,state]         (A)              {$s$};
\node[state]         (B) [right of=A] {$h$};

\path (A) edge [bend left] node {\TMtrans{A}{A}{\TMright}} (B)
(B) edge [loop above] node {\TMtrans{\TMblank}{\TMblank}{\TMright}} (B)
edge [bend left] node {\TMtrans{A}{A}{\TMright}} (A);
\end{tikzpicture}
\end{center}
\caption{\emph{सम} यंत्राची रूपे}
\ollabel{fig:variants}
\end{figure}
हे दोन आलेख दोन यंत्रे दर्शवतात: $M$ ची फितीची
वर्णमाला $\Sigma = \{\TMendtape,\TMblank,\TMstroke\}$
आणि अवस्थांचा संच $\{q_0,q_1\}$ आहे; तर $M'$ ची
वर्णमाला $\Sigma' = \{\TMendtape,\TMblank,A\}$ आणि
अवस्था $\{s,h\}$ आहेत. पण त्यांच्यातील सूचना
इतर बाबतीत सारख्याच आहेत: $n$ $\TMstroke$
चिन्हांची क्रमिका दिल्यावर $n$ सम असेल तेव्हाच
$M$ थांबते; आणि $n$ $A$ चिन्हांची क्रमिका दिल्यावर
$n$ सम असेल तेव्हाच $M'$ थांबते. आपण फक्त
$\TMstroke$ चे नाव~$A$, $q_0$ चे नाव~$s$
आणि $q_1$ चे नाव~$h$ असे बदलले आहे. हे
उदाहरण सर्वसाधारण करता येते: चिन्हे आणि अवस्था
अशा रीतीने पुनर्नामित केल्याने एक यंत्र दुसऱ्यात
रूपांतरित होत असेल, तर त्या ट्यूरिंग यंत्रांना
एकच मानता येते. खरे तर, ट्यूरिंग यंत्राची चिन्हे
आणि अवस्था या सरळ धन पूर्णांकच मानता येतात:
$\sigma_0$ ऐवजी~$1$, $\sigma_1$ ऐवजी~$2$
इत्यादी; $\TMendtape$ म्हणजे~$1$, $\TMblank$
म्हणजे~$2$ इत्यादी. अशा प्रकारे, \emph{सम}
यंत्र \olref{fig:standard-even} मध्ये दाखवलेल्या
यंत्रात रूपांतरित होते.
\begin{figure}
\[\begin{tikzpicture}[->,>=stealth',shorten >=1pt,auto,node distance=2.8cm,
  semithick]
\tikzstyle{every state}=[fill=none,draw=black,text=black]

\node[initial,state]         (A)              {$1$};
\node[state]         (B) [right of=A] {$2$};

\path (A) edge [bend left] node {\TMtrans{3}{3}{\TMright}} (B)
(B) edge [loop above] node {\TMtrans{2}{2}{\TMright}} (B)
edge [bend left] node {\TMtrans{3}{3}{\TMright}} (A);
\end{tikzpicture}
\]
\caption{एक मानक \emph{सम} यंत्र}
\ollabel{fig:standard-even}
\end{figure}
ज्या ट्यूरिंग यंत्राच्या अवस्था आणि चिन्हे धन पूर्णांक
असतात त्याला \emph{मानक} यंत्र म्हणता येईल; यापुढे
केवळ मानक यंत्रांचाच विचार करू.\footnote{“मानक
यंत्र” ही संज्ञा स्वतः मानक संज्ञा नाही.}

ट्यूरिंग यंत्रांचा संच !!{enumerable} म्हणजे गणनीय
आहे, हे दाखवण्याचे आपले उद्दिष्ट होते. वरील विचार
लक्षात घेतल्यास मानक ट्यूरिंग यंत्रांचा संच
!!{enumerable} म्हणजे गणनीय
आहे हे दाखवणे पुरेसे आहे. समजा, एक मानक ट्यूरिंग
यंत्र $M = \tuple{Q, \Sigma, q_0, \delta}$ दिले
आहे. धन पूर्णांकांच्या एका सांत क्रमिकेत त्याचे
वर्णन कसे करावे? प्रथम अवस्थांची संख्या, मग
अवस्था, चिन्हांची संख्या, चिन्हे आणि आरंभीची
अवस्था अशी यादी करू. (लक्षात घ्या, $M$ मानक
यंत्र असल्यामुळे ही सर्व धन पूर्णांक आहेत.)
मग~$\delta$ चे काय? संभाव्य आदाने, म्हणजे
$\tuple{q,\sigma}$ या जोड्या, यांचा संच सांत आहे;
कारण $Q$ आणि~$\Sigma$ सांत आहेत. म्हणून~$\delta$
मधील माहिती म्हणजे ज्या सर्व बाबतीत
$\delta(q,\sigma) = \tuple{q', \sigma', D}$ आहे,
अशा $\tuple{q, \sigma, q', \sigma', d}$
या सर्व $5$-घटकी क्रमितांची सांत यादी होय.
येथे $d$ ही दिशा~$D$ सांकेतिक
करणारी संख्या आहे (उदाहरणार्थ,
~$\TMleft$ साठी $1$, ~$\TMright$ साठी $2$
आणि ~$\TMstay$ साठी $3$).

अशा प्रकारे प्रत्येक मानक ट्यूरिंग यंत्राचे
धन पूर्णांकांच्या सांत यादीने, म्हणजे
$s_M \in (\PosInt)^*$ या क्रमिकेने, वर्णन
करता येते. उदाहरणार्थ, मानक \emph{सम}
यंत्र पुढील क्रमिकेने सांकेतिक केले जाते:
\[
2, \underbrace{1, 2}_Q, 3, \overbrace{1, 2, 3}^\Sigma, 1, \underbrace{1, 3, 2, 3, 2}_{\delta(1,3) = \tuple{2,3,R}}, 
\overbrace{2, 2, 2, 2, 2}^{\delta(2,2) = \tuple{2,2,R}},
\underbrace{2, 3, 1, 3, 2}_{\delta(2,3) = \tuple{1,3,R}}.
\]
\end{explain}

\begin{thm}
$\Nat$ पासून~$\Nat$ मध्ये जाणारी अशी फलने
आहेत जी ट्यूरिंग-संगणनक्षम नाहीत.
\end{thm}

\begin{proof}
धन पूर्णांकांच्या सांत क्रमिकांचा संच~$(\PosInt)^*$
!!{enumerable} म्हणजे गणनीय आहे, हे आपल्याला माहीत आहे
(\cref{sfr:siz:zigzag:prob:posint-star}). मानक ट्यूरिंग
यंत्रांच्या वर्णनांचा संच हा~$(\PosInt)^*$ चा उपसंच
असल्याने तोही गणनीय आहे. $\Nat$ पासून~$\Nat$
मध्ये जाणारे प्रत्येक ट्यूरिंग-संगणनक्षम फलन कोणत्या
ना कोणत्या (खरे तर अनेक) ट्यूरिंग यंत्रांनी संगणित
केले जाते. त्याच्या अवस्था आणि चिन्हांना धन
पूर्णांकांची नावे देऊन (विशेषतः $\TMendtape$ ला~$1$,
$\TMblank$ ला~$2$ आणि $\TMstroke$ ला~$3$)
प्रत्येक ट्यूरिंग-संगणनक्षम फलन कोणत्या ना कोणत्या
मानक ट्यूरिंग यंत्रानेही संगणित केले जाते, हे
दिसते. म्हणून $\Nat$ पासून~$\Nat$ मध्ये जाणाऱ्या
सर्व ट्यूरिंग-संगणनक्षम फलनांचा संचही गणनीय आहे.

दुसरीकडे, $\Nat$ पासून~$\Nat$ मध्ये जाणाऱ्या
सर्व फलनांचा संच !!{enumerable} म्हणजे गणनीय नाही
(\cref{sfr:siz:red:prob:nat-nat}). जर प्रत्येक
फलन एखाद्या ट्यूरिंग यंत्राने संगणित केले
जात असेल, तर ती फलने संगणित करणाऱ्या यंत्रांची
सर्व वर्णने यादीत दिल्याने सर्व फलनांचे प्रगणन
करता आले असते. त्यामुळे काही फलने
ट्यूरिंग-संगणनक्षम नाहीत.
\end{proof}

\begin{prob}
  अवस्था आणि वर्णमालेतील चिन्हांची स्वतंत्र
  स्पष्ट यादी न देता ट्यूरिंग यंत्रांचे वर्णन
  करण्याचा मार्ग तुम्हाला सुचतो का? “मानक”
  यंत्राची स्वतःची संकल्पना तुम्ही ठरवू शकता;
  पण तुमच्या नव्या अर्थाने प्रत्येक ट्यूरिंग
  यंत्राचे संगणन एखाद्या “मानक” यंत्राने
  कसे करता येईल, हेही स्पष्ट करा.
\end{prob}

\end{document}
